\section{The Technical Pipeline: From Waveform to (Something Like) Meaning}
\label{sec:pipeline}

If genuine translation were possible, what would the system look like?
Drawing on the cross-species decoding literature, we decompose the problem
into four stages---\emph{encoding}, \emph{structure discovery},
\emph{semantic grounding}, and \emph{interactive validation}---and assess
the maturity of each for the canine case. The stages are not a waterfall so
much as a loop: validation feeds back into grounding, which constrains
structure, which shapes the encoder.

\subsection{Stage 1: self-supervised acoustic encoding}
The modern starting point is a self-supervised representation trained on
large, unlabeled audio. AVES (Animal Vocalization Encoder based on
Self-Supervision) adapts the HuBERT/wav2vec paradigm to animal sound and
shows that representations learned without labels transfer to downstream
bioacoustic tasks~\cite{hagiwara2023aves}. NatureLM-audio extends this to an
\emph{audio--language} foundation model: trained on text--audio pairs
spanning bioacoustics, speech, and music, it can detect vocalizations,
classify rare species, and even caption or generate animal sounds from
natural-language prompts, demonstrating transfer from speech and music to
bioacoustics~\cite{robinson2024naturelm}. Google DeepMind's Perch line and
the global birdsong embeddings show that a single encoder can serve many
taxa~\cite{ghani2023perch}. Standardized benchmarks---BEANS (the Benchmark
of Animal Sounds) and BirdSet---give these encoders a common evaluation
surface~\cite{hagiwara2023beans}, while open general-purpose encoders such
as OpenBEATs push state of the art on a majority of BEANS
subtasks~\cite{openbeats2025}. For dogs, the implication is that a strong
\emph{Rung 1--2} encoder is available off the shelf; the bottleneck is not
representation quality but labels and semantics.

\subsection{Stage 2: sequence-structure discovery}
Given an encoder, the next question is whether the signal has
\emph{combinatorial structure}---discrete units that recombine, the way
phonemes form words. This is where the flagship projects have produced
their most cited results. Project CETI's analysis of 8{,}719 sperm-whale
codas revealed a ``phonetic alphabet'': rhythm, tempo, rubato, and
ornamentation combine contextually, expanding the repertoire of
distinguishable forms by nearly an order of magnitude~\cite{sharma2024whale}.
Crucially, the authors are explicit that this describes \emph{form}, not
\emph{meaning}: the paper answers what the units are and how they combine,
not what most of them say. Google's DolphinGemma takes the language-model
route directly---a roughly 400-million-parameter Gemma-based model with a
SoundStream tokenizer, trained on the Wild Dolphin Project's four-decade
Atlantic spotted dolphin corpus, that predicts the next sound in a sequence
exactly as an LLM predicts the next token~\cite{google2025dolphingemma}.
CETI's WhAM (2025) similarly learns to complete masked clicks and to render
whale-like codas from other audio. The Earth Species Project's 2026
zebra-finch model (ZF-AIM) goes further, detecting, predicting, and
\emph{generating} calls in real time and engaging birds in turn-taking
exchanges.

The lesson for dogs is double-edged. On one hand, the methods are
species-agnostic and could be applied to bark sequences. On the other, the
structure these models find is \emph{statistical}, and predicting the next
sound is not understanding it. Yossi Yovel and others have warned that such
models may cluster calls by acoustic similarity while missing meaning
entirely---grouping signals that sound alike but signify different
things. For a repertoire as graded and context-bound as the dog's, the
prior probability of finding a clean combinatorial alphabet is lower than
for whale codas or birdsong.

\subsection{Stage 3: semantic grounding}
Structure without semantics is syntax without a dictionary. Grounding ties
a signal to the world: to the eliciting context, the dog's subsequent
behavior, the objects in view, the social partner. This is the hardest
stage and the one where dogs are paradoxically \emph{advantaged}, because
they live in richly instrumented human homes and their behavior is legible
to us. The soundboard paradigm is effectively a grounding
device~\cite{bastos2024plos,bastos2024scirep}: the button press, the spoken
word, and the consequent action (door opens, toy appears) form a
triangulated, checkable mapping. The Wild Dolphin Project's CHAT system
does the analogous thing underwater, pairing a synthetic whistle with a
coveted object (sargassum, a scarf) so that the dolphin can request it by
name~\cite{google2025dolphingemma}. The general principle---\emph{ground
symbols through interaction and consequence}, not through passive
correlation---is the field's most important methodological insight, and it
is exactly what consumer bark-to-sentence apps lack.

\subsection{Stage 4: interactive validation (playback)}
The final stage closes the loop: does the animal treat the model's output
as meaningful? The gold standard is the \emph{playback experiment}---play a
recorded or synthesized signal and observe whether the receiver responds
appropriately. Comparative ethology has used playbacks for decades to
establish that, e.g., vervet alarm calls refer to specific predators. The
Earth Species Project's interactive finch work and the CHAT dolphin
protocol are modern, model-driven versions. For dogs, a validation loop
would look like: infer an intent, emit a candidate signal or action, and
test whether the dog's behavior confirms the inference---with the
experimenter blind to the model's guess to avoid Clever-Hans contamination.
No consumer product does this; it is the dividing line between science and
entertainment.

\subsection{Maturity summary}
Table~\ref{tab:pipeline} rates each stage for the canine case. The pattern
is clear: the \emph{front} of the pipeline (encoding, structure) is
borrowable from the broader bioacoustics boom and is relatively mature; the
\emph{back} (grounding, validation) is where canine-specific, falsifiable
work is needed and where almost none of the commercial effort is directed.
This asymmetry---strong models, weak semantics, absent validation---is the
technical signature of the whole field.

\begin{table}[htbp]
\centering
\caption{Maturity of the four decoding stages for the canine case.
``Borrowable'' means general bioacoustic methods transfer directly;
``canine-specific'' means dedicated dog data and protocols are required.}
\label{tab:pipeline}
\small
\begin{tabular}{@{}p{2.7cm}p{4.6cm}p{2.5cm}p{3.0cm}@{}}
\toprule
\textbf{Stage} & \textbf{Representative method} & \textbf{Maturity} & \textbf{Canine status} \\
\midrule
Encoding & AVES, NatureLM-audio, Perch~\cite{hagiwara2023aves,robinson2024naturelm,ghani2023perch} &
High & Borrowable; Rung 1--2 ready \\
Structure & CETI alphabet, DolphinGemma, WhAM~\cite{sharma2024whale,google2025dolphingemma} &
Medium & Statistical only; meaning open \\
Grounding & Soundboards, CHAT object-labeling~\cite{bastos2024plos,google2025dolphingemma} &
Low--medium & Needs interaction data \\
Validation & Playback / closed-loop response & Low & Absent in products \\
\bottomrule
\end{tabular}
\end{table}